% QMB_n13_hilbertraum.tex — Chapter 13: Hilbert Space Translation
% Substantive basis: Doc. 230 (May 2026), Doc. 231; own presentation
\chapter{Two Languages, One Content: FFGFT and Hilbert Space}
\label{ch:hilbert}

\section{The Claim}

The FFGFT describes quantum systems in the language of modes on a
torus. Standard quantum mechanics describes them in the language of
vectors in a Hilbert space. This chapter shows that both languages
carry the same content: Every calculation that is possible in one
is possible in the other, with the same result. The translation is
not an approximation and not a specialization, but a bijective
mapping of states, operators, gates, product spaces, and
measurement prescriptions.\status{B}

This is methodologically the most important statement of this book. It means
that the FFGFT does not replace quantum mechanics, but extends it:
Both formalisms yield, up to a $\xi$ correction of order
$10^{-5}$, the same measurable predictions. They differ in the
interpretation of some basic concepts --- Born rule, nonlocality, the status
of $\hbar$, the role of space-time, singularities. Anyone at home in
Hilbert space language can use the FFGFT without changing their
tools. The bridge is the translation table.

\section{Where We Are Geometrically}

Before the tables come, a picture of the geometry in which the translation
takes place. We will calculate with cylindrical coordinates, and the cylinder
is not the full FFGFT geometry, but a convenient reduction.

\paragraph{The full geometry is four-dimensional.} The space on which
all statements of the FFGFT live is the torus $T^4 = T^3\times S^1$: three
spatial winding directions and one temporal. The wave function
of a mode $\varphi_{n,l,j}$ is a function on this torus. From
its geometry follow the masses, the fine structure constant, the
cosmological quantities --- everything that distinguishes the FFGFT from pure
quantum mechanics.

\paragraph{A single qubit needs less.} A two-level system
at fixed time, without mass dynamics, does not activate most of these
winding directions. Two rotational degrees of freedom remain: one that
encodes where between $|0\rangle$ and $|1\rangle$ we are (this becomes $z$),
and one for the phase (this becomes $\theta$). The natural space for this is a 2-torus $T^2 = S^1\times S^1$.

\paragraph{From the 2-torus to the cylinder.} If the major radius $R$ of the
2-torus is large compared to the tube radius, the geometry looks locally like a
cylinder $\mathbb{R}\times S^1$. One of the two loops is
cut open: Two circles become one line and one circle. In
cylindrical coordinates $(z,r,\theta)$, $z$ is therefore a linear axis
from $-1$ to $+1$, no longer cyclic.

\paragraph{From the cylinder to the Bloch sphere.} If one pulls the cylinder together at the top
and bottom into one point each, a sphere $S^2$ arises: the
Bloch sphere of standard quantum mechanics. It is the second,
further reduction.

\begin{center}
\small\renewcommand{\arraystretch}{1.3}
\begin{tabular}{p{2.9cm}cp{3.6cm}}
\toprule
\textbf{Geometry} & \textbf{Dim.} & \textbf{Use} \\
\midrule
4-torus $T^4$ & 4 & full FFGFT: masses, $\alpha$, cosmology \\
3-torus $T^3$ & 3 & mode space at fixed time \\
2-torus $T^2$ & 2 & single qubit without reduction \\
Cylinder $\mathbb{R}\times S^1$ & 2 & limit $R\to\infty$, local approximation \\
Bloch sphere $S^2$ & 2 & poles pulled together; standard QM \\
\bottomrule
\end{tabular}
\end{center}

\paragraph{What the reductions cost.} Each level loses something. From the
4-torus to the 3-torus, the time evolution is lost and with it every
statement about masses. From the 3-torus to the 2-torus, one spatial
winding is lost and with it the connection to the third lepton generation
and to many spin structures. From the 2-torus to the cylinder, one
compact loop is lost --- that is the weakest reduction, and
exactly here sits the fractal correction factor
$K_{\mathrm{frak}} = 1-100\xi\approx0{,}9867$, which appears in the
bit-flip rotations. $K_{\mathrm{frak}}$ is the cylinder's memory of its
torus origin.

The translation tables of this chapter apply to the qubit sector,
i.e., to the bottom three rows. Standard QM works on the
Bloch sphere, the FFGFT on the cylinder or the 2-torus. The
translation between these levels is lossless, provided one carries the
small geometric corrections ($K_{\mathrm{frak}}$, $\Delta$CHSH).
The 4-torus lies one level higher and provides what the
Hilbert space cannot provide: the values of the input parameters.

\begin{laierspur}
One can draw a map in different ways: as a globe,
as an unrolled cylinder, as a flat sheet. All show the same
cities in the same places. But only the globe shows that one
arrives back in the west from the east. The Bloch sphere of
quantum mechanics is the flat sheet; the torus of the FFGFT is the globe.
For navigation between two cities, the sheet suffices. For the
question of why the world is as large as it is, one needs the globe.
\end{laierspur}

\section{States}

In standard QM, a qubit is a vector in $\mathbb{C}^2$:
\begin{equation}
  |\psi\rangle = \alpha|0\rangle + \beta|1\rangle,\qquad
  |\alpha|^2+|\beta|^2 = 1,\qquad \alpha,\beta\in\mathbb{C}.
\end{equation}
With global phase freedom, this is a point on the Bloch sphere
(Hopf fibration $S^3\to S^2$).

In the FFGFT, the same qubit is a point $(z,r,\theta)$ on the
cylinder:
\begin{equation}
  z\in[-1,1],\quad r\in[0,1],\quad \theta\in[0,2\pi),\qquad z^2+r^2 = 1.
\end{equation}
$z$ is the projection onto the computation axis, $r$ the radial
superposition amplitude, $\theta$ the phase.

The bridge between the two representations:
\begin{equation}
  \boxed{\;
  \alpha = \sqrt{\tfrac{1+z}{2}}\;e^{\,i\theta/2},\qquad
  \beta = \sqrt{\tfrac{1-z}{2}}\;e^{-i\theta/2}.\;}
  \label{eq:state-bridge}
\end{equation}
Conversely:
\begin{equation}
  z = |\alpha|^2-|\beta|^2,\qquad
  r = 2|\alpha\beta|,\qquad
  \theta = \arg\alpha-\arg\beta.
\end{equation}
The consistency is checked in four lines: The normalization
$|\alpha|^2+|\beta|^2 = \frac{1+z}{2}+\frac{1-z}{2} = 1$ holds; the
cylinder condition $z^2+r^2 = (|\alpha|^2-|\beta|^2)^2+4|\alpha|^2|\beta|^2 = 1$
holds; $|0\rangle$ corresponds to $(1,0,\ast)$ and $|1\rangle$ corresponds to
$(-1,0,\ast)$, with arbitrary phase at the poles.

The conceptual difference lies not in the numbers. In QM,
$|\alpha|^2$ is a probability. In the FFGFT, $r^2 = 1-z^2$
is an energy density of the radial configuration. Both quantities are numerically equal for large statistics; conceptually, the FFGFT is a
deterministic geometric model whose statistical predictions
coincide with the Born rule.

\section{Operators}

The three Pauli matrices generate all single-qubit operations:
\begin{equation}
  \sigma_x = \begin{pmatrix}0&1\\1&0\end{pmatrix},\quad
  \sigma_y = \begin{pmatrix}0&-i\\i&0\end{pmatrix},\quad
  \sigma_z = \begin{pmatrix}1&0\\0&-1\end{pmatrix},
\end{equation}
with the algebra $\sigma_i\sigma_j = \delta_{ij}\mathbb{1}+i\epsilon_{ijk}\sigma_k$.

On the cylinder, three geometric rotations correspond to them:
$\sigma_z$ is the rotation about the cylinder axis, i.e., the phase rotation
$\theta\to\theta+\pi$; $\sigma_x$ is the rotation in the $(z,r)$ plane by
the angle $\pi$, the bit flip; $\sigma_y$ is the rotation orthogonal to it, which
combines bit flip and phase rotation.

\begin{center}
\small\renewcommand{\arraystretch}{1.5}
\begin{tabular}{cp{2.6cm}p{3.6cm}}
\toprule
\textbf{Operator} & \textbf{Hilbert Space} & \textbf{FFGFT Geometry} \\
\midrule
$\sigma_z$ & $\mathrm{diag}(1,-1)$ & phase rotation $\theta\to\theta+\pi$ \\
$\sigma_x$ & $\begin{pmatrix}0&1\\1&0\end{pmatrix}$ & $(z,r)$ rotation by $\pi K_{\mathrm{frak}}$ \\
$\sigma_y$ & $\begin{pmatrix}0&-i\\i&0\end{pmatrix}$ & rotation with phase flip \\
$H$ & $\frac{1}{\sqrt2}\begin{pmatrix}1&1\\1&-1\end{pmatrix}$ & swap $z\leftrightarrow r$, $\theta\to\theta+\pi/2$ \\
$T$ & $\mathrm{diag}(1,e^{i\pi/4})$ & phase rotation $\theta\to\theta+\pi/4$ \\
\bottomrule
\end{tabular}
\end{center}

The only FFGFT-specific deviation is the factor
$K_{\mathrm{frak}}\approx0{,}9867$ in the $\sigma_x$ rotation. It follows from
the fractal space-time dimension $D_f = 2{,}999867$ and is a testable
prediction: All $\pi$ gates deviate systematically by
$\Delta\alpha\approx0{,}013\,\%$ from the ideal value.\status{S} In the
Hilbert space formalism, this would be an ad hoc correction that would have to be built in as hardware noise; in the FFGFT it follows from
the geometry.

\section{Gates}

Every unitary single-qubit gate $U\in U(2)$ is a Bloch rotation:
\begin{equation}
  U = e^{i\alpha}R_{\hat n}(\phi)
    = e^{i\alpha}\bigl(\cos\tfrac{\phi}{2}\,\mathbb{1}
      - i\sin\tfrac{\phi}{2}\,\hat n\cdot\vec\sigma\bigr).
\end{equation}
In the FFGFT, this is the rotation of the point $(z,r,\theta)$ about the axis
$\hat n$ by the angle $\phi$. The translation is exact.

The CNOT gate
\begin{equation}
  \mathrm{CNOT} =
  \begin{pmatrix}1&0&0&0\\0&1&0&0\\0&0&0&1\\0&0&1&0\end{pmatrix}
\end{equation}
translates as a prescription: If $z_A = -1$ (control qubit at
$|1\rangle$), perform the bit flip on qubit B. This is a geometric
conditioning --- the rotation of one cylinder point depends on the
$z$ value of the other. Topologically, this is a linking of the two
cylinders, and that is the entanglement.

The set $\{H,T,\mathrm{CNOT}\}$ is universal for quantum mechanics according to the Solovay-Kitaev theorem.\status{E} Since each of these gates is
representable in the FFGFT as a geometric operation, the FFGFT is
fully universal for quantum computation: There is no QM statement
that cannot be formulated in FFGFT language.\status{B}

\section{Product Spaces and Entanglement}

For $n$ qubits, the Hilbert space is
$\mathcal{H}_n = \bigotimes_{i=1}^n\mathbb{C}^2$ with dimension $2^n$,
and a general state reads
\begin{equation}
  |\Psi\rangle = \sum_{x\in\{0,1\}^n}\alpha_x|x\rangle,\qquad
  \sum_x|\alpha_x|^2 = 1.
\end{equation}
In the FFGFT, $n$ qubits correspond to $n$ mode configurations
$(z_i,r_i,\theta_i)$ on a common torus geometry. Entanglement
is a topological connection: Entangled qubits are knots of the
mode configuration that cannot be reduced to product form.
The Bell state $|\Phi^+\rangle = \frac{1}{\sqrt2}(|00\rangle+|11\rangle)$
is a topologically connected cylinder configuration in which the
individual $(z_i,r_i,\theta_i)$ are not independent.

Both formalisms reproduce the exponential growth $\dim = 2^n$:
in the Hilbert space through the tensor product, in the FFGFT through $n$
independent mode configurations with topological connections. The
number of real parameters per state is $2\cdot2^n-2$ in both cases.

The CHSH prediction: Standard QM gives $|\langle\mathrm{CHSH}\rangle|
\le2\sqrt2$. The FFGFT predicts a small, geometrically motivated deviation,
$\Delta\mathrm{CHSH}\approx\xi/(2\pi)\approx10^{-5}$ per
measurement.\status{S} It lies below the current hardware noise threshold;
the distinction from the cumulative quantity from Doc.~022 is explained in
Chapter~\ref{ch:chsh_exp}.

\section{Measurement}

A measurement in the computational basis yields in QM the value $0$ with
probability $|\alpha|^2$ and the value $1$ with probability
$|\beta|^2$; the state collapses to $|0\rangle$ or $|1\rangle$.

In the FFGFT, measurement is not a probabilistic collapse, but a
geometric interaction that drives the point $(z,r,\theta)$ to the
nearest stable extremum $z = \pm1$. This is a
deterministic movement in mode space --- because of the unknown phase
$\theta$ and smallest variations of the initial conditions, statistically
indistinguishable from the behavior of the Born rule. For large statistics,
the measurement distribution converges:
\begin{equation}
  P_{\mathrm{FFGFT}}(z = +1) = \tfrac{1+z}{2} = |\alpha|^2 = P_{\mathrm{QM}}(0).
\end{equation}
The formalisms yield identical statistical predictions.\status{B}
The dispute over local realism resolves in the FFGFT geometrically,
because entanglement is represented as a topological connection and not as a
nonlocal correlation (Chapter~\ref{ch:bell_ffgft}).

\section{Time Evolution}

The time evolution of QM is the Schrödinger equation
$i\hbar\,\partial_t|\psi\rangle = H|\psi\rangle$. In the FFGFT, it is the
deterministic movement of the mode configuration $\varphi_{n,l,j}$ on
$T^3\times\mathbb{R}$. The mode equation reduces in the
low-energy limit $E\ll E_\xi$ to the Schrödinger equation:
\begin{equation}
  i\hbar\,\frac{\partial\varphi_{n,l,j}}{\partial t}
  \;\xrightarrow{\;E\ll E_\xi\;}\;
  -\frac{\hbar^2}{2m}\nabla^2\varphi_{n,l,j}+V\varphi_{n,l,j}.
\end{equation}
At higher energies, $\xi$ corrections occur, which in the
Hilbert space formalism appear as modifications of the Hamilton operator.

\section{What Is Translatable and What Is Not}

All Hilbert space statements are mappable in FFGFT: States, operators
and gates (with or without $K_{\mathrm{frak}}$), tensor products,
entanglement, Bell and GHZ states, measurement, unitary evolution and
all algorithms. Conversely, every $(z,r,\theta)$ configuration corresponds uniquely
to an $(\alpha,\beta)$ vector, every geometric rotation to a
unitary operator, every topological connection to an entangled
vector. Every calculation that is performable in one formalism is
performable in the other --- with identical result.

The translatability holds for statements about quantum systems. It does
not hold for statements about the input values of quantum mechanics:
\begin{itemize}
  \item The parameter $\xi = 4/30000$ is an external
        input value in the Hilbert space. In the FFGFT it follows from the geometry
        (Chapter~\ref{ch:grundlage}).
  \item The fermion masses appear in the Hilbert space as Yukawa couplings.
        In the FFGFT they follow from $m_i = r_i\,\xi^{p_i}\,v$.
  \item The fine structure constant is derivable in the FFGFT, input in the
        Hilbert space.
  \item $K_{\mathrm{frak}}$ appears in FFGFT gates as
        geometric deviation; in the Hilbert space it would have to be modeled as noise.
\end{itemize}
That is the structural asymmetry: The Hilbert space has all tools
for QM calculations, but no explanation for the input values. The FFGFT
provides the tools and the input values from a single geometry.
In this sense, quantum mechanics is a subset of the FFGFT: the
FFGFT, restricted to the question of how systems behave, without the
question of why the constants have these values.

The FFGFT does not claim to be the only correct representation. It is a representation that additionally provides the values of the
fundamental constants. Anyone who does not need this added value
can ignore it; standard QM remains valid.

\section{Operational Equivalence, Interpretive Divergence}

The complete translatability must not be read as ontological identity.
Operationally, FFGFT and standard QM agree in every measurable result up to $\xi$ corrections. Interpretively,
they diverge at five points where the standard reading makes statements
that are read in the FFGFT as artifacts of the chosen carrier assumptions:

\begin{enumerate}
  \item \textbf{Born rule as irreducible chance.} Standard reading:
        $|\psi|^2$ is ontological probability, the measurement outcome
        fundamentally random. FFGFT: $|\psi|^2$ is the correct distribution
        of observations, but arises from a deterministic pole transition dynamics. The stochasticity sits in the measuring apparatus,
        not in the natural process.
  \item \textbf{Nonlocality as a property of nature.} Standard reading
        of the Bell correlations: action at a distance in $\mathbb{R}^3$. FFGFT:
        topological connection on $T^4$, geometric proximity in a
        closed geometry.
  \item \textbf{$\hbar$ as primitive.} Standard reading: fundamental
        constant of nature. FFGFT: SI conversion factor between energy and
        time measure, which follows from the closure condition $S = h$ on $T^4$.
        Numerically identical, ontologically derived.
  \item \textbf{Space-time as stage.} Standard reading of the theory of
        relativity: given manifold. FFGFT: derived from the
        $T^4$ identification geometry. Questions like
        \enquote{what was before the Big Bang} are artifacts of the
        $\mathbb{R}^3\!+\!t$ assumption.
  \item \textbf{Singularities as realities.} Standard reading: Black
        hole and Big Bang as true singularities. FFGFT: Since $T^4$
        is closed, they are endpoints of the carrier model, not of the
        physics. Outside the singularity region, the predictions
        agree.
\end{enumerate}

These five points are not empirical claims against standard
QM. They are interpretation differences that become visible
as soon as one asks about the ontological carrier. Anyone who only calculates with the numbers
can ignore them. Anyone who wants to know why the numbers are
these specific ones will find them to be testable structural statements.

Within today's measurement precision --- NISQ level $\sim10^{-2}$,
optical precision experiments $\sim10^{-12}$ in Bell tests --- both
yield the same predictions. The $\xi$ correction $\sim10^{-5}$ is
not resolvable with current hardware; this has been empirically tested
on IBM Heron~r2 (Chapter~\ref{ch:chsh_exp}).

\section{Field-Theoretic Extension}

The translation can be extended from the qubit sector to quantum field states.
The Fock space of standard QFT --- the direct sum of all
$n$-particle spaces --- corresponds to the mode spectrum of the carrier, where
the occupation numbers appear as winding numbers. Creation and
annihilation operators are transitions between neighboring
winding states of the same mode; commutators follow from the
periodicity. This generates no new predictions in the perturbative
regime, but restructures the construction of the Standard Model geometrically:
The particle species are mode classes of the lattice, not independent
fields with their own coupling each.\status{S}

\begin{keypoint}[What This Chapter Holds]
States, operators, gates, entanglement, measurement, and time evolution
are bijectively translatable between FFGFT and Hilbert space. The only
formal deviation is $K_{\mathrm{frak}}$ in the $\pi$ gates; the
only statistical one is $\Delta$CHSH. The added value of the FFGFT lies not
in different calculation results, but in the fact that the input values of
quantum mechanics follow from the same geometry.
\end{keypoint}

\quelle{Doc.~230 (Translatability, May 2026), Doc.~231 (field-theoretic extension), Doc.~147 §1.2, Doc.~175 §3--4, Doc.~202 §22}